\section{Введение}
Методы Франка–Вульфа (FW) решают линейные подзадачи вместо вычисления
проекций. Обратный поиск позволяет применять квадратичные модели шага
без заранее заданной константы гладкости, а шаги удаления улучшают
продвижение на многогранниках~\cite{braun,pedregosa,lacostejaggi}.
Однако глобальная граница кривизны может терять существенную информацию,
когда кривизна убывает вместе со значением целевой функции.
Мы рассматриваем аффинное ограничение гессиана через отклонение потерь
от оптимума, изученное в~\cite{warmup}, применительно к выпуклой задаче
с ограничениями. Оно также является случаем показателя один в семействе
гладкости, зависящей от потерь, из~\cite{liuwarmup}.
Более общие ограничения гессиана через значение функции
ранее возникли в теориях самоограничивающейся
регулярности~\cite{kale,caoselfbound}; само условие не является
вкладом настоящей работы.

Наш вопрос количественный: что гарантирует обычный обратный поиск
для наблюдаемого правила остановки при таком условии, с учётом
начальной оценки и всех пробных вычислений? Теорема~\ref{thm:fw}
отвечает на него, разделяя вклады $H_0/\varepsilon$ и
$H_1\log(1/\varepsilon)$. Основной инструмент — потенциал оценок,
которые остаются слишком большими при переходе сертификата на более
низкие двоичные уровни. Он не меняет алгоритм и не ограничивает LMO
множеством подуровня. Теорема~\ref{thm:afw} переносит результат с
помощью стандартного подсчёта удаляемых атомов. При геометрическом
неравенстве для зазора теорема~\ref{thm:geometric} даёт логарифмическую
зависимость от точности с явной начальной ценой $H_1$.
Теорема~\ref{lem:bridge} устанавливает оптимальную поточечную огибающую
FW-зазора для этого класса и одномерные функции, на которых она
достигается. Утверждение~\ref{prop:zeroface} характеризует случай
$H_0=0$: всё множество решений является экспонированной гранью с общим
градиентом. Теорема~\ref{thm:zerofinite} затем усиливает приближённую
остановку до конечного достижения точного оптимума методом AFW с
обратным поиском, с явной оценкой, не зависящей от запрошенной точности.
Используются известные механизмы идентификации~\cite{bomze,garber};
новый алгоритм AFW и общая теория идентификации не заявляются.
Приведены полные доказательства, граничные примеры и контрпримеры.

\paragraph{Связанные исследования.}
Гладкость, зависящая от функции, и двухрежимные оценки градиентных
методов предшествуют этой работе~\cite{warmup,lobanovkoloskova,lobanov}.
Выгузов и Стонякин~\cite{vyguzov} анализируют FW при гладкости,
зависящей от нормы градиента. На компактном выпуклом множестве
$f-f_Q^*\le D\norm{\nabla f}$ переводит наше условие в соответствующую
границу гессиана, но глобальная оценка максимального градиента не
сохраняет изучаемое здесь разделение зависимостей.
Pedregosa и соавторы~\cite[теоремы 1, 3 и 4]{pedregosa} уже дают
подсчёт проб обратного поиска, двойственную гарантию Фенхеля с
усреднённым градиентом и линейную сходимость AFW с принятыми локальными
оценками кривизны. Наш максимум скалярных нижних оценок отличается от
этой двойственной величины; их значения нельзя просто отождествить.
Наш результат специализирует локальную кривизну к порогу сертификата
и явно учитывает запаздывание. Приоритет обратного поиска, двойственных
сертификатов, линейной сходимости AFW и телескопического тождества для
отклонённых проб не заявляется.
Лучшие накопленные касательные границы и вычислимые зазоры остановки
также известны ранее~\cite[метод 1]{freundgrigas}. Оптимистический FW и
прямые-двойственные короткие шаги явно оптимизируют вычислимые зазоры
при глобальной гладкости~\cite[утверждение 4.1]{beyondshort}; они не
оценивают анализируемое здесь унаследованное правило деления и
умножения оценки на два.

К другим связанным подходам относятся FW для обобщённо
самосогласованных функций, адаптивные шаги Брегмана и адаптивный
слайдинг условного градиента~\cite{selfconcordant,bregmanfw,adcgs}.
Оценки расширенной кривизны и схемы с рестартами при условиях оценки
ошибки уже дают двухрежимные и геометрические скорости~\cite{wirth,restartfw};
линейная сходимость по функции и FW-зазору за пределами глобальной
гладкости градиента также известна для логарифмически однородных
барьеров~\cite{zhao}. Автоматически обусловленный FW служит практическим
вариантом сравнения с иной предварительной обработкой и учётом
оракулов~\cite{autoconditioned}. Инициализация по разности градиентов с
обратным поиском~\cite{khademi} и условный градиент без линейного
поиска~\cite{yagishita} дают другие способы адаптации к параметрам.
ALFCG оценивает гладкость без линейного поиска, но сохраняет условия
глобальной покомпонентной или средней гладкости и стохастический
критерий стационарности~\cite[раздел 3]{alfcg}.
Недавний Newton--FW с демпфированием использует информацию гессиана и
предполагает его положительную определённость~\cite[раздел 1]{newtonfw};
это не охватывает наш оптимум с нулевой базовой кривизной, где гессиан
обращается в нуль.
Приложение~\ref{app:comparison} уточняет используемые различия.
Прямая подстановка глобальной гладкости обычно слабее, но это не
означает, что эквивалентное следствие невозможно вывести из другого
адаптивного анализа.

Экспериментальная часть отвечает на три вопроса, существенных для
интерпретации оценок: помогает ли инициализация по сертификату независимо
от смены стартовой точки; устраняют ли шаги удаления препятствия
активной грани; означает ли отсутствие остановки оставшуюся ошибку по
функции или консервативность сертификата. Бюджеты оракулов учитываются
раздельно; неуспешные запуски сохраняются.
